% Q1: Series RLC Natural Response — SPDT switch moves from position a to b at t=0
\documentclass[border=10pt]{standalone}
\usepackage[american voltages, american currents]{circuitikz}
\usepackage{amsmath}
\renewcommand{\familydefault}{\sfdefault}

\begin{document}
\begin{circuitikz}[
    line width=0.8pt,
    every node/.style={font=\sffamily}
]

% Voltage source (left, vertical, positive on top)
\draw (0,4) to[V, v=$V_s$] (0,0);

% Wire from Vs+ to SPDT position a
\draw (0,4) -- (1.2,4);

% ---- SPDT Switch SW1 ----
% Position a: (1.2, 4)   — connects to Vs+
% Position b: (1.2, 1.5) — connects to bottom rail (closes RLC loop)
% Common:    (3.0, 3.5)  — connects to RLC chain

% Terminal dots
\fill (1.2,4) circle (2pt);
\fill (1.2,1.5) circle (2pt);
\fill (3.0,3.5) circle (2pt);

% Switch arm from common toward position a (t < 0 state, solid)
\draw[thick] (3.0,3.5) -- (1.35,3.95);

% Labels for positions
\node[above, font=\sffamily\small] at (1.2,4.1) {\textit{a}};
\node[left, font=\sffamily\small] at (1.05,1.5) {\textit{b}};
\node[above right, font=\sffamily\itshape\small] at (3.0,3.55) {SW1};
\node[font=\sffamily\footnotesize] at (2.6,4.3) {$t = 0$};

% Wire from position b down to bottom rail
\draw (1.2,1.5) -- (1.2,0);
\fill (1.2,0) circle (2pt);

% Wire from common up to top rail, then right to RLC
\draw (3.0,3.5) -- (3.0,4) -- (6,4);

% Series R-L-C going down on the right
\draw (6,4) to[R, l_=$R$] (6,2.7)
            to[L, l_=$L$] (6,1.3)
            to[C, l_=$C$, v^=$v_C$] (6,0);

% Bottom return rail
\draw (6,0) -- (0,0);

% Current arrow on top rail
\draw[->, >=stealth, thick] (3.5,4.3) -- (5.5,4.3)
    node[midway, above]{$i(t)$};

\end{circuitikz}
\end{document}
